% !TeX root = ./00-main.tex
% ============================================================================
% Chapter 1 -- Introduction  [FINAL, trimmed pass]
%
% Cuts applied throughout: removed restated clauses (sentences that repeat the
% preceding sentence's claim in different words), dead lead-ins to citations
% ("Recent studies have begun to..."), and redundant scene-setting ahead of
% concrete evidence. No claims, citations, or hedges were removed -- only
% wording. Section-by-section rationale was reviewed with the author before
% reassembly; see prior turn for the paragraph-by-paragraph diff and reasoning.
%
% ONE STRUCTURAL NOTE NOT YET ACTED ON: the Kiemen2026-dz "tens of microns"
% resolution fact is now made in Sec 1.1 para 4 (as neighbor evidence), Sec 1.2
% para 3 (as the core resolution argument), and echoed in Sec 1.3's closing
% paragraph. That's three appearances of one fact. Left as-is pending your call
% on whether to state it once and cross-reference, or keep the repetition
% because each occurrence serves a different local argument.
% ============================================================================

\chapter{Introduction}
\label{chap:intro}

\section{Premalignant progression is a tissue process}
\label{sec:intro-premalignancy}

Two recent mechanistic studies show epithelial and microenvironmental states
evolving together with explicit direction and feedback, rather than as parallel
change. In mouse lung, \emph{KRAS}-mutant alveolar type II cells secrete
amphiregulin, activate EGFR signaling in adjacent fibroblasts, and induce a
fibrotic injury-like program; the reprogrammed fibroblasts then alter alveolar
macrophage state and reinforce epithelial plasticity, and disrupting the
circuit prevents tumor initiation \autocite{Cardoso2026-kp}. In mouse pancreas,
oncogenic and tumor-suppressive programs converge within a progenitor-like
population at the benign-to-malignant transition; suppressing p53 in that
population permits its expansion, epithelial--mesenchymal reprogramming, and
formation of an immune-privileged niche \autocite{Reyes2026-yr}. In both
systems the epithelium does not simply prompt stromal change and receive it
back passively: it constructs a microenvironment that then acts back on the
epithelial state that built it.

Lung adenocarcinoma provides the primary system for this thesis. Its
recognized precursor sequence extends from atypical adenomatous hyperplasia
(AAH) through adenocarcinoma in situ (AIS) and minimally invasive
adenocarcinoma to invasive disease, and multiregion, single-cell, spatial,
genomic, and epigenomic studies document coordinated remodeling of epithelial,
immune, stromal, and vascular compartments across it
\autocite{Sinjab2021-in,Zhu2022-hz,Haga2023-xc,Chen2025-zo}. \emph{KRT8}-positive
alveolar intermediate states define a plastic epithelial axis associated with
injury, regeneration, and early transformation \autocite{Wang2023-hg}, and
ablating a high-plasticity state in early neoplasia prevents the
benign-to-malignant transition \autocite{Chan2026-ov}. In the spatial
multi-omics cohort analyzed here, alveolar progenitor-like states localize
within proinflammatory niches enriched for \emph{IL1B}-high macrophages; these
niches are prevalent in precursor lesions, become less frequent in established
adenocarcinoma, promote progenitor growth, and can be reduced by inflammatory
intervention \autocite{Peng2025-xu} --- a stage-specific distribution that
places the niche within early progression rather than as a consequence of
established cancer.

Pancreatic intraepithelial neoplasia provides a complementary ductal precursor
system with substantially different tissue ecology. PanINs are widespread,
multifocal, and genetically heterogeneous in grossly normal adult pancreata,
yet only a minority give rise to invasive disease
\autocite{Braxton2024-sv,Kiemen2024-xk}. Their progression occurs within a
fibroblast-rich, spatially organized environment marked by ductal remodeling,
acinar disruption, local inflammation, immune reorganization, and stromal
restructuring \autocite{Carpenter2023-qs,Bell2024-kn,Lyman2025-qg,Kiemen2026-dz},
with fibroblast state, inflammatory signaling, and immune architecture
changing across PanIN grade and progression toward pancreatic ductal
adenocarcinoma \autocite{Bell2024-kn,Lyman2025-qg}. Where LUAD is organized
around epithelial--macrophage inflammatory niches, PanIN progression is
organized more strongly around fibroblast, matrix, and lesion-associated
immune remodeling --- a stringent test of whether a context-effect framework
transfers across distinct precursor ecologies rather than recovering one
tissue-specific architecture.

PanIN also reveals a consequential separation between epithelial and
ecological progression. In donor pancreata, PanIN epithelial cells already
show substantial transcriptional similarity to malignant epithelial cells,
while their surrounding fibroblast and macrophage states remain distinct from
those of pancreatic cancer \autocite{Carpenter2023-qs}. Spatial analyses extend
this by placing the PanIN epithelial compartment on a continuum with invasive
cancer while its microenvironment stays sharply separated and reorganizes only
with malignant progression \autocite{Elhossiny2026-pq}; failure to acquire a
cancer-permissive stromal state has been proposed as one reason most PanINs
remain indolent \autocite{Elhossiny2026-pq}. If progression depends on such a
state emerging, the relevant ecological information may reside not in an
epithelial receiver's immediate neighborhood but in a broader lesion- or
specimen-level field. Distinguishing receiver-local context from
specimen-scale ecology is therefore not a modeling decision but a biological
question about the spatial organization of premalignant progression.


\section{What spatial omics measures, and what it does not}
\label{sec:intro-measurement}

Single-cell and single-nucleus RNA sequencing resolve epithelial and
microenvironmental states at high molecular resolution; spatial
transcriptomics situates those states within tissue architecture. Together
they supply the three components context-associated progression requires: a
regulatory representation of the epithelial receiver, a representation of its
surrounding tissue, and an ordered pathological axis along which state change
can be inferred. They do not measure the transition connecting these
components directly; it must be reconstructed from the structure of the data.

The first constraint is spatial resolution. Spot-based spatial
transcriptomics aggregates expression across multicellular regions, so it
does not directly observe a single epithelial cell together with a uniquely
defined microenvironment. Reference mapping and deconvolution instead
estimate the cellular composition and continuous cell-state structure
contributing to each measurement, using methods such as cell2location,
Tangram, RCTD, TACCO, and DestVI
\autocite{Kleshchevnikov2022-ki,Biancalani2021-bt,Cable2022-tb,Mages2023-wj,Lopez2022-ps}.
The resulting representation depends jointly on the reference population, its
annotation granularity, the mapping model's assumptions, and the platform's
physical resolution --- choices that determine which epithelial states can
serve as receivers and which ecological features can be assigned to their
surroundings, and so define the model's biological state space rather than a
neutral preprocessing layer.

This limitation is biologically consequential because tissue ecology varies
across spatial scales. Cellular-resolution three-dimensional histology of more
than one thousand human pancreatic precancers reveals marked heterogeneity in
perilesional inflammation, with immune hot and cold spots alternating over
tens of microns \autocite{Kiemen2026-dz}. Spot-based platforms integrate
expression over an effective length scale of tens to roughly one hundred
microns, so spatial variation below that scale is not directly observed ---
and reference-based inference can estimate a spot's cellular constituents but
cannot reconstruct subspot organization that was never measured. Any context
representation derived from these data therefore inherits a resolution
boundary that added model complexity cannot remove.

The second constraint is temporal. Premalignant spatial cohorts are
cross-sectional because profiling consumes the tissue: successive stages are
represented by different cells, lesions, and patients, producing
stage-indexed population distributions rather than longitudinal trajectories
of individual cells. No receiver is observed both before and after
progression. The transition between stages must instead be inferred through a
correspondence model relating unpaired source and target populations
\autocite{Bunne2024-pl,Weinreb2018-jq} --- a model-based quantity, not an
empirical path recovered from repeated measurement, whose interpretation
depends on the molecular representation, the assumed transition geometry, and
the structure imposed on correspondence. The objective of this thesis follows
from these two constraints: not merely to fit stage-associated variation, but
to determine which components of an inferred transition remain identifiable
under explicit falsification.

\section{The unresolved estimand}
\label{sec:intro-gap}

Spatial transcriptomics has generated a rapidly expanding methodological
literature for characterizing the organization of tissue ecosystems.
Neighborhood and niche-detection methods identify recurrent cellular
configurations and spatially organized tissue states; latent-space approaches
resolve location-specific molecular variation associated with
tumor--microenvironment interactions \autocite{Deshpande2023-zw}; and
neighborhood-augmented representations such as BANKSY incorporate local
molecular context directly into the representation of each spatial
observation \autocite{Singhal2024-rv}. In parallel, optimal transport and
flow-matching methods align molecular populations across conditions and
estimate the transition fields connecting them
\autocite{Cuturi2013-li,Schiebinger2019-lm,Lipman2022-ns,Klein2025-sr,Richter2026-ib}.

These methodological families resolve complementary aspects of tissue
organization but leave a specific quantity undefined. Spatial niche methods
identify where particular cellular configurations occur and which molecular
programs are associated with them. Transport methods estimate how populations
are related across stages or conditions. Neither, by itself, estimates the
context-associated component of progression for a fixed epithelial receiver:
the difference in its predicted transition when tissue context is included,
after holding its initial regulatory state and pathological stage edge
constant. Nor do these approaches determine whether such a quantity is
identifiable from cross-sectional spatial measurements.

This distinction introduces three inferential failure modes. First, spatial
enrichment does not imply an effect on transition: a niche may become more
prevalent at a later stage because its constituent cell states expand
together, even if it contributes no information about how an individual
receiver changes. Second, improved prediction after adding neighborhood
features does not establish that neighborhood \emph{identity} is informative
--- a context-conditioned model may outperform a self-only one simply from
additional features, parameters, or specimen-level ecology correlated with
stage. Third, the correspondence used to define the transition can absorb the
ecological signal before the context effect is estimated: if source--target
matching is itself determined by the context features under evaluation, the
resulting estimate cannot distinguish context-associated displacement from
context-conditioned correspondence.

A defensible context-effect analysis must address all three problems
simultaneously. It requires an explicitly defined within-receiver estimand, a
correspondence model constructed so that the effect of interest is not
encoded upstream, and a falsification test that destroys receiver--context
correspondence while preserving the marginal ecological structure supplied to
the model. The decisive comparison is therefore not simply between a model
with context and a model without it. It is between the observed context
assignment and an appropriately matched context assignment whose distribution
is retained but whose receiver identity has been removed.

The same logic makes spatial scale part of the estimand rather than an
implementation detail. A model supplied with receiver-local context can
exploit information carried by the immediate neighborhood, by a broader
tissue region, or by specimen-level ecological composition without
distinguishing among them. Neighborhood-scale conditioning is therefore a
hypothesis about where progression-relevant information resides, not a
property guaranteed by the input representation. This assumption is placed
under tension by the sub-spot ecological variation already noted in
Section~\ref{sec:intro-measurement} and by evidence that
progression-permissive stromal states may emerge as lesion- or
specimen-level fields \autocite{Elhossiny2026-pq}. Whether the recoverable
signal is receiver-specific, specimen-scale, or inaccessible below the
resolution of the platform is consequently an empirical question. Resolving
it requires a framework capable not only of detecting an apparent context
effect, but also of concluding that receiver-local context carries no
identifiable information beyond the ecological distribution from which it
was drawn.

\section{Correspondence and dynamics from cross-sectional populations}
\label{sec:intro-dynamics}

Relating stage-indexed populations requires a model of correspondence between
unpaired molecular distributions. Pseudotime estimates relative ordering
within an observed state manifold, and RNA velocity estimates short-horizon
direction from transcriptional kinetics; neither addresses the case in which
source and target populations arise from different tissue sections, lesions,
and patients \autocite{Richter2026-ib,Weinreb2018-jq}.

Optimal transport supplies a principled framework for that case. A transport
coupling assigns probability mass between source and target states under a
specified transition geometry, and entropic regularization makes estimation
tractable in high-dimensional molecular data \autocite{Cuturi2013-li}.
Waddington-OT established the framework for destructive single-cell time
courses, and subsequent work extended it across complete temporal series
\autocite{Schiebinger2019-lm,Lavenant2023-qv}. The resulting coupling depends
on the molecular representation, the transport cost, the regularization, and
the marginal structure --- choices that determine which earlier states are
associated with which later states, and which therefore require
justification rather than default settings.

Continuous dynamical models convert discrete correspondences into fields of
state change. Neural ordinary differential equations parameterize a state
derivative through a learned vector field \autocite{Chen_undated-rf}; flow
matching estimates such fields directly from interpolated source--target
pairs \autocite{Tong2020-ou,Huguet2022-pz,Lipman2022-ns}. Optimal-transport
conditional flow matching combines the two, using a transport coupling to
define training pairs and a continuous field to connect their distributions
\autocite{Klein2025-sr,Richter2026-ib} --- so pathological stage becomes an
ordered coordinate along which population-level regulatory states transform.

Spatial conditioning has begun to enter this setting. NicheFlow models the
evolution of structured cellular microenvironments across spatial time
courses, and MIOFlow~2.0 conditions cellular trajectories on neighbourhood
composition and intercellular signalling \autocite{Sakalyan2025-ia,Sun2026-xk}.
What such conditioning introduces, and what the present work takes as its
central concern, is a question of scale: context may describe an ecological
state shared across an entire specimen, a regional tissue configuration, or
the immediate neighbourhood of an individual receiver, and a predictive model
can exploit any of these without distinguishing among them. Specimen-level
ecology, receiver-local context, and donor composition are three different
sources of progression-associated information, and separating them is a
design requirement rather than a post-hoc analysis.

\section{StageBridge and the conditional context residual}
\label{sec:intro-estimand}

StageBridge estimates progression-aligned regulatory transport from
separately encoded representations of epithelial receiver state and tissue
context. Each focal epithelial receiver is represented by its regulatory
state, while ecological information enters the dynamical model through a
dedicated context branch. This architectural separation allows intrinsic
epithelial state, specimen-level ecology, and receiver-local context to be
evaluated as distinct sources of progression-associated information.

The central estimand is defined by contrasting two trajectories initialized
from the same epithelial receiver. Integrating the fitted dynamical field
with the context branch active produces one endpoint; integrating from the
same initial state with the context branch inactive produces another. Their
difference defines the conditional context residual,

\begin{equation}
R_{\mathrm{cond},i}
=
z_i^{\mathrm{full}}(1)
-
z_i^{\mathrm{self}}(1),
\label{eq:intro-context-residual}
\end{equation}

where $z_i^{\mathrm{full}}(1)$ and $z_i^{\mathrm{self}}(1)$ denote the
endpoints reached from the same initial receiver state under the full and
self-only dynamics, respectively. Because the two trajectories share the same
initial condition, stage edge, and fitted model parameters, the residual is a
within-receiver contrast rather than a comparison between independently
fitted models. It quantifies the endpoint shift associated with activating
the fitted context branch while holding the receiver's starting regulatory
state fixed. Chapter~\ref{chap:methods} specifies the estimand formally, and
Appendix~\ref{chap:appendix-b} provides its derivation.

The interpretation of this residual depends on three design choices. First,
the dynamics are evaluated in a named regulatory space composed of curated
pathway and transcription-factor activities rather than in an unconstrained
latent embedding. Conditional residuals can therefore be reported as signed,
program-resolved quantities; PROGENy and decoupleR provide the corresponding
regulatory coordinates \autocite{Schubert2018-ph,Badia-I-Mompel2022-te}.
Second, the progression field is parameterized as a structured universal
differential equation that combines specified regulatory terms with a
flexible learned residual component \autocite{Rackauckas2020-mc,Philipps2025-jl}.
This construction preserves the distinction between the explicitly modeled
and learned contributions to the field and permits their relative magnitudes
to be quantified. Third, source--target correspondence is estimated within
coarse ecological strata rather than optimized directly on the full
receiver-local context. The transport coupling therefore constrains
plausible correspondences without encoding, upstream, the same ecological
information whose downstream contribution the model is intended to estimate.

The evaluation framework is designed to distinguish predictive fit from
identifiability. Donor-held-out prediction, specimen-average context
substitution, matched context shuffling, architectural ablation, and a
frozen synthetic recovery benchmark jointly test whether an estimated context
association generalizes, the spatial scale at which its information is
carried, and whether the estimator recovers a known effect under controlled
conditions. The decisive comparison is between the observed
receiver--context assignment and a matched assignment whose context identity
has been destroyed but whose marginal ecological structure is preserved ---
the only contrast that can separate receiver-specific information from added
model capacity and the ecological distribution the context was drawn from.

\section{Objective and hypotheses}
\label{sec:intro-objective}

The objective of this thesis is to determine whether tissue context contains
information about progression-aligned epithelial regulatory transport beyond
that contained in the initial epithelial state, to identify the spatial scale
at which any such information is recoverable, and to resolve the regulatory
programs through which it is expressed. Lung adenocarcinoma precursors
provide the primary system for model development and biological discovery.
Pancreatic intraepithelial neoplasia provides a secondary evaluation of
estimator behavior and transferability within a biologically distinct
precursor ecology.

Three hypotheses were stated conceptually here and formalized in
Chapter~\ref{chap:methods}, Section~\ref{sec:hypotheses}. Chapter~\ref{chap:results}
evaluates each hypothesis against its prespecified criterion.

\paragraph{Hypothesis 1: a reproducible conditional context residual exists.}
After conditioning on the epithelial receiver state and ordered pathological
transition, tissue context is associated with a reproducible, program-resolved
component of progression-aligned regulatory displacement. The corresponding
prediction is that specific named regulatory programs exhibit conditional
context residuals whose direction and magnitude remain stable across
donor-held-out resampling.

\paragraph{Hypothesis 2: progression-associated context information is
receiver-local.}
The spatial context assigned to an individual epithelial receiver contains
progression-associated information beyond both specimen-level ecology and the
receiver's intrinsic regulatory state. The corresponding prediction is that a
model supplied with each receiver's observed local context will predict the
target-stage distribution more accurately than models supplied with
specimen-average context, no context, or matched context whose
receiver-specific correspondence has been disrupted while its marginal
distribution is preserved. The final comparison constitutes the stringent
test: failure to distinguish observed from shuffled context would indicate
that receiver identity contributes no detectable endpoint-predictive
information beyond the ecological distribution from which the context was
sampled.

\paragraph{Hypothesis 3: the estimand exhibits transferable behavior across
precursor systems.}
When applied without modification to a biologically distinct precursor
epithelium, the same estimand should produce a coherent result rather than an
artifact of the primary cohort. It should identify interpretable regulatory
programs where the secondary cohort supports stable estimation and return no
stable programs where it does not. This analysis tests the behavior and
portability of the estimator; it does not constitute independent validation
of the primary lung findings, and its interpretation is bounded by the size
and structure of the secondary cohort.

The hypotheses address related but non-equivalent properties of the
framework: the reproducibility of the estimated residual, the spatial
localization of the information it captures, and the behavior of the
estimand across biological systems. The thesis is organized around their
separate evaluation rather than presuming that all three resolve
concordantly.

\section{Organization of the thesis}
\label{sec:intro-organization}

Chapter~\ref{chap:methods} defines the biological cohorts and analysis units;
constructs the epithelial receiver and tissue-context representations;
specifies the conditional-stratified transport coupling, structured dynamical
model, and conditional context residual; and formalizes the evaluation
framework and hypothesis criteria.

Chapter~\ref{chap:results} presents the empirical analysis in four parts. It
first evaluates whether the estimand is recoverable and whether its behavior
is an artifact of the regulatory representation. It then characterizes the
program-resolved residuals identified across the two lung progression edges,
examines estimator behavior in pancreatic intraepithelial neoplasia, and
compares the inferred programs with complementary analyses of the same
biological systems. The chapter next determines whether the recoverable
ecological information is carried by receiver-local context or by broader
specimen-level structure. It concludes by resolving each prespecified
hypothesis against the resulting inferential boundary.

Chapter~\ref{chap:discussion} integrates the biological and methodological
findings. It interprets the regulatory programs nominated by the conditional
residual, develops the spatial-scale result as the central inferential
contribution, defines the limitations imposed by cross-sectional sampling and
platform resolution, and identifies the experimental and computational
developments required to distinguish local signaling effects from broader
ecological organization.

Appendix~\ref{chap:appendix-a} derives entropic optimal transport and the
Sinkhorn iteration. Appendix~\ref{chap:appendix-b} derives the dynamical
training objective and conditional context residual and provides the
complete training procedure. Appendix~\ref{chap:appendix-c} contains the
dataset inventory, regulatory feature definitions, fold assignments,
hyperparameters, complete falsification ladders, program-stability atlas,
supplementary analyses, and reproducibility manifest.


% ============================================================================
% CITATION KEYS USED IN THIS CHAPTER
%
% Verified present in the LUAD/PanIN .bib exports:
%   Braxton2024-sv  Kiemen2024-xk   Teixeira2019-yb  Mazzilli2025-hi
%   Pedro2025-kz    Anderson2026-zg Cardoso2026-kp   Reyes2026-yr
%   Kadara2016-dt   Sinjab2021-in   Haga2023-xc      Chen2025-zo
%   Wang2023-hg     Chan2026-ov     Peng2025-xu      Bell2024-kn
%   Lyman2025-qg    Kiemen2026-dz   Elhossiny2026-pq
%
% Carried over from the previous Ch1 draft (confirm these resolve):
%   Kleshchevnikov2022-ki  Biancalani2021-bt  Cable2022-tb   Mages2023-wj
%   Lopez2022-ps           Bunne2024-pl       Weinreb2018-jq Richter2026-ib
%   Cuturi2013-li          Schiebinger2019-lm Lavenant2023-qv Chen_undated-rf
%   Tong2020-ou            Huguet2022-pz      Lipman2022-ns  Klein2025-sr
%   Singhal2024-rv         Deshpande2023-zw   Sakalyan2025-ia Sun2026-xk
%   Schubert2018-ph        Badia-I-Mompel2022-te Rackauckas2020-mc
%   Philipps2025-jl        Carpenter2023-qs
%
% (NEEDS KEY) -- Sec 1.3 cites the niche-detection literature generically. The
% claim is stronger with specific references. Add and cite:
%   - NicheCompass / niche quantification        (Birk et al., Nat Genet 2025)
%   - COVET/ENVI, covariance environment          (Haviv et al., Nat Biotech 2025)
%   - Niche-DE, niche differential expression     (Mason et al. 2023)
%   - single-cell niche identification            (Qian et al., Nat Commun 2025)
% These four are in the thesis proposal bibliography but were not in the two
% .bib exports; pull them into the main library.
%
% NOTE: \autocite{Chen_undated-rf} has an "undated" key -- fix the entry.
% ============================================================================